\chapter{Dialogs}

Standard dialogs are one of the highest-value parts of the \api{crystal-qt6} library because they unlock native-feeling file selection, confirmation flows, color picking, font selection, text input, and progress reporting. A desktop application can feel dramatically more complete once these flows use the platform's expected dialogs.

This chapter is intentionally screenshot-heavy. The figures are build-safe placeholders until the screenshots are captured, but each placeholder names the image file that should eventually be stored under \api{docs/book/images}.

\section{Dialog Gallery}

The maintained gallery at \api{examples/dialog_gallery.cr} opens each supported dialog from buttons, menu actions, and toolbar actions. It is the recommended source for documentation screenshots because it keeps all dialog flows in one small, reproducible application.

\begin{lstlisting}[style=crystal]
make example-dialogs
\end{lstlisting}

\screenshotplaceholder{docs/book/images/dialogs-main-window.png}{The dialog gallery main window with buttons, menus, toolbar actions, and result labels.}

\section{Convenience Helpers And Explicit Instances}

Most dialog APIs offer a convenience helper for the common modal flow. The helper creates the dialog, configures it, runs it, returns the result, and releases the native dialog.

\begin{lstlisting}[style=crystal]
if name = Qt6::InputDialog.get_text(
  main,
  title: "Rename Layer",
  label: "Layer name",
  value: "Terrain"
)
  status_bar.show_message("Renamed to #{name}", 1500)
end
\end{lstlisting}

Use an explicit dialog instance when you need more configuration, signals, or post-acceptance inspection.

\begin{lstlisting}[style=crystal]
dialog = Qt6::FileDialog.new(main, Dir.current, "Maps (*.map *.json);;All Files (*)")
dialog.window_title = "Open Map"
dialog.accept_mode = Qt6::FileDialogAcceptMode::Open
dialog.file_mode = Qt6::FileDialogFileMode::ExistingFile

begin
  if dialog.exec == Qt6::DialogCode::Accepted
    status_bar.show_message(dialog.selected_file)
  end
ensure
  dialog.release
end
\end{lstlisting}

The choice is mostly about how much control the workflow needs. Prefer a helper when the dialog has one job, the result is enough, and no caller needs to observe intermediate signals. Prefer an explicit instance when you need to preselect values, tune modes and options, connect callbacks, inspect multiple results, or keep ownership and release behavior visible in the surrounding code.

\begin{itemize}
\item Use convenience helpers for ordinary modal questions such as ``choose one file,'' ``pick a color,'' or ``enter one number.''
\item Use explicit instances for customized file modes, progress cancellation, font/color signals, custom button boxes, and flows that need several selected values.
\item Parent explicit dialogs to the window that launched them. For short-lived unparented dialogs, release them in an \api{ensure} block.
\item In tests, explicit instances are usually easier to configure, show, drive with timers, and release predictably.
\end{itemize}

\section{Message Boxes}

\api{MessageBox} is for confirmation, warnings, errors, and short informational messages. The convenience helpers match common message-box intent: \api{information}, \api{warning}, \api{critical}, and \api{question}.

\begin{lstlisting}[style=crystal]
answer = Qt6::MessageBox.question(
  main,
  title: "Save Changes?",
  text: "The current map has unsaved changes.",
  informative_text: "Save before opening another file?",
  buttons: Qt6::MessageBoxButton::Yes |
           Qt6::MessageBoxButton::No |
           Qt6::MessageBoxButton::Cancel
)

case answer
when Qt6::MessageBoxButton::Yes
  save_project
when Qt6::MessageBoxButton::No
  open_next_project
else
  status_bar.show_message("Open canceled", 1500)
end
\end{lstlisting}

\smallscreenshotplaceholder{docs/book/images/dialogs-message-question.png}{A question message box with Yes, No, and Cancel buttons.}

\begin{lstlisting}[style=crystal]
Qt6::MessageBox.information(
  main,
  title: "Export Complete",
  text: "The preview image was written successfully."
)

Qt6::MessageBox.critical(
  main,
  title: "Export Failed",
  text: "The preview image could not be written.",
  informative_text: "Check that the destination folder is writable."
)
\end{lstlisting}

\begin{figure}[htbp]
  \begin{lstlisting}[style=crystal]
Qt6::MessageBox.warning(
  main,
  title: "Layer Warning",
  text: "The selected layer is locked.",
  buttons: Qt6::MessageBoxButton::Ok |
           Qt6::MessageBoxButton::Cancel
)
  \end{lstlisting}
  \centering
  \screenshotimage{docs/book/images/dialogs-message-warning.png}{width=0.42\textwidth,height=0.24\textheight,keepaspectratio}%
  \caption{A warning message box for a recoverable problem.}
\end{figure}

\clearpage

For noninteractive testing or deeper configuration, create \api{MessageBox} directly, set \api{text}, \api{informative_text}, \api{icon}, and \api{standard_buttons}, then call \api{show_modal}.

\section{File Dialogs}

\api{FileDialog} covers open, multi-open, save, and explicitly configured file picker flows. File dialogs should use a starting directory and a useful name filter whenever possible.

\screenshotplaceholder{docs/book/images/dialogs-file-open.png}{An open-file dialog with Crystal source and all-file filters.}
\screenshotplaceholder{docs/book/images/dialogs-file-save.png}{A save-file dialog configured for project output.}

\begin{lstlisting}[style=crystal]
path = Qt6::FileDialog.get_open_file_name(
  main,
  "Open Project",
  Dir.current,
  "Project files (*.json *.map);;All Files (*)"
)

if path
  status_bar.show_message("Opened #{File.basename(path)}", 1500)
end
\end{lstlisting}

\begin{lstlisting}[style=crystal]
paths = Qt6::FileDialog.get_open_file_names(
  main,
  "Import Assets",
  Dir.current,
  "Images (*.png *.jpg *.jpeg);;All Files (*)"
)

status_bar.show_message("Selected #{paths.size} asset files", 1500)
\end{lstlisting}

\begin{lstlisting}[style=crystal]
path = Qt6::FileDialog.get_save_file_name(
  main,
  "Export Preview",
  File.join(Dir.current, "preview.png"),
  "PNG Images (*.png);;All Files (*)"
)

export_preview(path) if path
\end{lstlisting}

Explicit instances expose \api{accept_mode}, \api{file_mode}, \api{directory}, \api{name_filter}, \api{select_file}, \api{selected_file}, and \api{selected_files}. Use them when a file workflow needs to be configured before the dialog is shown.

\begin{lstlisting}[style=crystal]
dialog = Qt6::FileDialog.new(main, Dir.current)
dialog.window_title = "Choose Asset Folder"
dialog.file_mode = Qt6::FileDialogFileMode::Directory

begin
  if dialog.exec == Qt6::DialogCode::Accepted
    import_folder(dialog.selected_file)
  end
ensure
  dialog.release
end
\end{lstlisting}

\section{Color Dialogs}

\api{ColorDialog} returns a \api{Color}. It is useful for theme colors, layer colors, brush colors, annotation colors, and any workflow where the user needs a visual picker.

\begin{lstlisting}[style=crystal]
accent = Qt6::Color.new(52, 112, 176)

if color = Qt6::ColorDialog.get_color(
  main,
  current_color: accent,
  title: "Pick Accent Color",
  show_alpha_channel: true
)
  accent = color
  status_bar.show_message(
    "Accent #{color.red},#{color.green},#{color.blue},#{color.alpha}",
    1500
  )
end
\end{lstlisting}

\smallscreenshotplaceholder{docs/book/images/dialogs-color-dialog.png}{A color dialog initialized with the current accent color.}
\smallscreenshotplaceholder{docs/book/images/dialogs-color-alpha.png}{A color dialog with alpha-channel editing enabled.}

\api{ColorDialog#native_dialog=} can disable the platform-native picker. That is mainly useful for deterministic tests, where the widget-backed dialog behaves like an ordinary Qt dialog.

Native color, font, and file dialogs are the right default for user-facing applications because they match platform expectations. The widget-backed versions are useful when screenshots, tests, or tutorial examples need a predictable Qt window instead of a system sheet or platform-owned picker. If a chapter screenshot needs to teach the binding API rather than the operating system's picker, prefer the widget-backed version and mention that real applications can leave native dialogs enabled.

\section{Font Dialogs}

\api{FontDialog} returns a \api{QFont}. It supports initial fonts, selected fonts, option flags, native-dialog control, and signals for highlighted or selected fonts.

\smallscreenshotplaceholder{docs/book/images/dialogs-font-dialog.png}{A font dialog initialized from the current label font.}
\smallscreenshotplaceholder{docs/book/images/dialogs-font-monospaced.png}{A font dialog filtered toward monospaced fonts.}

\begin{lstlisting}[style=crystal]
current_font = Qt6::QFont.new("Sans Serif", 12)

if font = Qt6::FontDialog.get_font(main, current_font, "Choose Label Font")
  current_font = font
  status_bar.show_message("#{font.family}, #{font.point_size} pt", 1500)
end
\end{lstlisting}

Use font-dialog options when the workflow cares about font category:

\begin{lstlisting}[style=crystal]
font = Qt6::FontDialog.get_font(
  main,
  Qt6::QFont.new("Monospace", 12),
  "Choose Code Font",
  Qt6::FontDialogOption::MonospacedFonts
)
\end{lstlisting}

\section{Input Dialogs}

\api{InputDialog} is a compact way to request one value from the user. The convenience helpers cover text, integers, floating-point values, and item choices.

\begin{lstlisting}[style=crystal]
name = Qt6::InputDialog.get_text(
  main,
  title: "Rename Layer",
  label: "Layer name",
  value: "Terrain"
)
\end{lstlisting}

\begin{lstlisting}[style=crystal]
opacity = Qt6::InputDialog.get_int(
  main,
  title: "Layer Opacity",
  label: "Opacity percent",
  value: 72,
  minimum: 0,
  maximum: 100
)
\end{lstlisting}

\begin{lstlisting}[style=crystal]
kind = Qt6::InputDialog.get_item(
  main,
  title: "Layer Kind",
  label: "Kind",
  items: ["Terrain", "Roads", "Units", "Annotations"],
  current: 1,
  editable: false
)
\end{lstlisting}

\begin{lstlisting}[style=crystal]
scale = Qt6::InputDialog.get_double(
  main,
  title: "Export Scale",
  label: "Scale factor",
  value: 1.0,
  minimum: 0.25,
  maximum: 4.0
)
\end{lstlisting}

\begin{lstlisting}[style=crystal]
tag = Qt6::InputDialog.get_item(
  main,
  title: "Layer Tag",
  label: "Tag",
  items: ["draft", "review", "final"],
  current: 0,
  editable: true
)
\end{lstlisting}

\inputdialogscreenshots

Use explicit \api{InputDialog} instances when you need to change the input mode, inspect range values, or configure combo-box editability directly.

\section{Progress Dialogs}

\api{ProgressDialog} reports long-running work and gives the user a chance to cancel. It has a label, cancel-button text, range, current value, auto-close behavior, auto-reset behavior, minimum display duration, and a canceled signal.

\screenshotplaceholder{docs/book/images/dialogs-progress-dialog.png}{A progress dialog updating from a timer.}
\screenshotplaceholder{docs/book/images/dialogs-progress-canceled.png}{A progress dialog after user cancellation.}

\begin{lstlisting}[style=crystal]
progress = Qt6::ProgressDialog.new(
  main,
  "Generating preview assets...",
  "Cancel",
  0,
  100
)
progress.window_title = "Export"
progress.minimum_duration = 0
progress.auto_close = true
progress.auto_reset = true
progress.value = 0

timer = Qt6::QTimer.new(main)
value = 0

progress.on_canceled do
  timer.stop
  status_bar.show_message("Export canceled", 1500)
end

timer.on_timeout do
  value += 10
  value = 100 if value > 100
  progress.value = value

  if value >= 100
    timer.stop
    status_bar.show_message("Export complete", 1500)
  end
end

progress.show
timer.start(100)
\end{lstlisting}

For real long-running work, make sure the GUI thread can continue processing events. A progress dialog that never repaints or never receives cancel input is usually a sign that too much work is blocking the event loop.

\section{Splash Screens}

\api{SplashScreen} is a lightweight startup surface. Use it only while the application is preparing its real main window; once the window can be shown, call \api{finish}.

\begin{lstlisting}[style=crystal]
pixmap = Qt6::QPixmap.new("assets/splash.png")
splash = Qt6::SplashScreen.new(pixmap)
splash.show
splash.show_message("Loading project...")
app.process_events

main = build_main_window
main.show
splash.finish(main)
\end{lstlisting}

Splash screens should stay brief. If the user might need to cancel, inspect details, or wait for a longer task, prefer a \api{ProgressDialog} instead.

\section{Custom Dialogs And Button Boxes}

Some workflows need a custom dialog body with standard buttons. \api{Dialog} provides accepted and rejected callbacks. \api{DialogButtonBox} provides standard button layouts and accept/reject signals.

\screenshotplaceholder{docs/book/images/dialogs-custom-settings.png}{A custom settings dialog with a form body and standard buttons.}

\begin{lstlisting}[style=crystal]
dialog = Qt6::Dialog.new(main)
dialog.window_title = "Layer Settings"

name = Qt6::LineEdit.new("Terrain")
visible = Qt6::CheckBox.new("Visible")
buttons = Qt6::DialogButtonBox.new(
  Qt6::DialogButtonBoxStandardButton::Ok |
  Qt6::DialogButtonBoxStandardButton::Cancel,
  dialog
)

dialog.vbox do |column|
  form_host = Qt6::Widget.new(dialog)
  form_host.form do |form|
    form.add_row("Name", name)
    form.add_row(visible)
  end

  column << form_host
  column << buttons
end

buttons.on_accepted { dialog.accept }
buttons.on_rejected { dialog.reject }
\end{lstlisting}

Custom dialogs should still follow the same ownership rules as other widgets: parent child controls to the dialog, keep the body focused, and let the button box handle standard accept/reject controls.

\section{Dialog Troubleshooting}

Most dialog problems are event-loop, ownership, or platform issues. These checks usually find the problem quickly:

\begin{itemize}
\item If a dialog never appears, make sure a \api{Qt6.application} exists and the dialog is launched from the GUI thread.
\item If a modal dialog appears behind another window, pass the launching window as the parent and avoid using an unparented top-level dialog for application-local questions.
\item If screenshots differ across macOS, Linux, and Windows, check whether the dialog is native. Native dialogs are allowed to look and behave like the host platform.
\item If a progress dialog shows but does not repaint, move long work into smaller chunks and let the event loop process between updates.
\item If a scripted capture misses a native dialog on macOS, run the gallery in a foreground desktop session and recapture manually when needed.
\end{itemize}

\section{Screenshot Checklist}

When replacing placeholders with real screenshots, capture the gallery in a consistent theme and window size. The first pass should include these files:

\begin{itemize}
\item \api{dialogs-main-window.png}
\item \api{dialogs-message-question.png}
\item \api{dialogs-message-warning.png}
\item \api{dialogs-file-open.png}
\item \api{dialogs-file-save.png}
\item \api{dialogs-color-dialog.png}
\item \api{dialogs-color-alpha.png}
\item \api{dialogs-font-dialog.png}
\item \api{dialogs-font-monospaced.png}
\item \api{dialogs-input-text.png}
\item \api{dialogs-input-int.png}
\item \api{dialogs-input-choice.png}
\item \api{dialogs-progress-dialog.png}
\item \api{dialogs-progress-canceled.png}
\item \api{dialogs-custom-settings.png}
\end{itemize}

The book Makefile tracks \api{docs/book/images/*.png}, so replacing a screenshot and rerunning \api{make docs-book} rebuilds the PDF. If you add a new image format or store generated screenshots elsewhere, update \api{docs/book/Makefile} so image changes remain visible to the build.

\section{Current Coverage}

The dialog surface includes \api{MessageBox}, \api{FileDialog}, \api{ColorDialog}, \api{FontDialog}, \api{InputDialog}, \api{ProgressDialog}, \api{SplashScreen}, \api{DialogButtonBox}, and base \api{Dialog} accepted/rejected callbacks.
